% Image credits for the photographs and AI illustrations of Book 4
% (University Biology, Year 2). Machine-readable license records live in
% images/book4/CREDITS.md; keep the two files in sync.
\chapter*{Créditos de las imágenes}

% \sloppy must be in force when the paragraph is BROKEN, hence the \par before
% \endgroup -- without it the group closes first and the tolerance is restored.
\begingroup\sloppy

Los dibujos de este libro son originales. Las fotografías se utilizan
bajo sus respectivas licencias libres, con agradecimiento a sus
autores:

\begin{itemize}
  \item Capítulo~\ref{ch:b2:microbial-metabolism}: fotografía de Sergei
        Winogradsky, década de 1890 --- dominio público (Wikimedia Commons).
  \item Capítulo~\ref{ch:b2:genome-diversification}: fotografía de
        Barbara McClintock en su laboratorio, Smithsonian Institution /
        Science Service, 1947 --- dominio público (Wikimedia Commons).
  \item Capítulo~\ref{ch:b2:blood-circulation}: retrato de William
        Harvey y lámina de las venas del antebrazo de su
        \emph{Exercitatio anatomica de motu cordis} (1628), Wellcome
        Collection --- CC~BY~4.0 (Wikimedia Commons); micrografía
        electrónica de barrido de un glóbulo rojo, una plaqueta y un
        linfocito, Electron Microscopy Facility, National Cancer
        Institute --- dominio público (Wikimedia Commons).
  \item Capítulo~\ref{ch:b2:speciation}: lámina de John Gould con los
        pinzones de Darwin, del \emph{Journal of Researches} (1845) ---
        dominio público (Wikimedia Commons).
  \item Capítulo~\ref{ch:b2:biogeochemical-cycles}: el observatorio del
        Mauna Loa (Forrest M. Mims III, 2006) y la toma de una muestra en
        matraz en ese lugar (comandante John Bortniak, 1982), NOAA Photo
        Library --- dominio público (Wikimedia Commons).
\end{itemize}

Los enlaces completos a las fuentes y las direcciones de las licencias
de cada fotografía figuran en \texttt{images/\allowbreak book4/\allowbreak CREDITS.md},
en el repositorio del libro.

\bigskip
Junto a las fotografías y los dibujos originales, unas setenta
figuras repartidas por los capítulos del volumen son ilustraciones
generadas por inteligencia artificial, producidas con la generación de
imágenes de OpenAI a partir de indicaciones escritas por los editores
del libro, y revisadas para comprobar su exactitud biológica antes de
incluirlas. La indicación completa de cada imagen generada está
registrada en \texttt{images/\allowbreak book4/\allowbreak ai/\allowbreak PROMPTS.md}, en el repositorio.
\par\endgroup
\clearpage
